\section{Affine combinations}
An affine combination is a linear combination whose coefficients sum to $1$. They may be negative: nonnegativity is an additional condition for a \emph{convex} combination, studied later. The affine hull $\operatorname{aff}S$ is the set of all affine combinations from $S$.

For nonempty $S=\{v_1,\ldots,v_k\}$,
\[
\operatorname{aff}S=v_1+\Span\{v_2-v_1,\ldots,v_k-v_1\}.
\]
An equivalent computational method appends a coordinate $1$ to each point. Solving $[\widetilde v_1\ \cdots\ \widetilde v_k]c=\widetilde p$ then enforces both the point equation and $\sum c_j=1$. An affine set is closed under affine combinations; it may be empty under the usual closure convention. A nonempty affine set is a translated subspace, or flat.

\practice{8.1.P1}{Three noncollinear points fill the plane affinely}
The edge differences from $v_1$ are $v_2-v_1=(-2,2)$ and $v_3-v_1=(2,1)$. Their determinant is $(-2)1-2\cdot2=-6\ne0$, so they span $\R^2$. Hence the affine hull of the three points is all of $\R^2$, and it must contain $p$.
\begin{center}\setlength{\unitlength}{6mm}
\begin{picture}(8.00,6.00)
\put(0.2,1.500){\vector(1,0){7.700}}
\put(2.500,0.2){\vector(0,1){5.700}}
\put(0.500,1.500){\line(0,1){.12}}\put(0.500,1.200){\makebox(0,0){\tiny -2}}
\put(1.500,1.500){\line(0,1){.12}}\put(1.500,1.200){\makebox(0,0){\tiny -1}}
\put(2.500,1.500){\line(0,1){.12}}\put(2.500,1.200){\makebox(0,0){\tiny 0}}
\put(3.500,1.500){\line(0,1){.12}}\put(3.500,1.200){\makebox(0,0){\tiny 1}}
\put(4.500,1.500){\line(0,1){.12}}\put(4.500,1.200){\makebox(0,0){\tiny 2}}
\put(5.500,1.500){\line(0,1){.12}}\put(5.500,1.200){\makebox(0,0){\tiny 3}}
\put(6.500,1.500){\line(0,1){.12}}\put(6.500,1.200){\makebox(0,0){\tiny 4}}
\put(7.500,1.500){\line(0,1){.12}}\put(7.500,1.200){\makebox(0,0){\tiny 5}}
\put(2.500,0.500){\line(1,0){.12}}\put(2.200,0.500){\makebox(0,0){\tiny -1}}
\put(2.500,2.500){\line(1,0){.12}}\put(2.200,2.500){\makebox(0,0){\tiny 1}}
\put(2.500,3.500){\line(1,0){.12}}\put(2.200,3.500){\makebox(0,0){\tiny 2}}
\put(2.500,4.500){\line(1,0){.12}}\put(2.200,4.500){\makebox(0,0){\tiny 3}}
\put(2.500,5.500){\line(1,0){.12}}\put(2.200,5.500){\makebox(0,0){\tiny 4}}
\linethickness{.9pt}
\polyline(3.500,1.500)(1.500,3.500)(5.500,2.500)(3.500,1.500)
\put(3.500,1.500){\circle*{.12}}\put(3.620,1.680){\scriptsize $v_1$}
\put(1.500,3.500){\circle*{.12}}\put(1.620,3.680){\scriptsize $v_2$}
\put(5.500,2.500){\circle*{.12}}\put(5.620,2.680){\scriptsize $v_3$}
\put(6.500,4.500){\circle*{.12}}\put(6.620,4.680){\scriptsize $p$}
\end{picture}\end{center}
For explicit weights solve $p-v_1=c_2(v_2-v_1)+c_3(v_3-v_1)$:
\[
-2c_2+2c_3=3,\qquad2c_2+c_3=3.
\]
Adding suitable multiples gives $c_2=1/2$, $c_3=2$, and then $c_1=1-c_2-c_3=-3/2$. Thus $p=-\tfrac32v_1+\tfrac12v_2+2v_3$. These weights sum to $1$ and reproduce $(4,3)$; a negative weight is permitted for an affine combination.

\exercise{8.1.1}{Require the weights to sum to one}
To require both the point equation and the coefficient-sum condition, solve
\[Hc=\widetilde y,\qquad H=\begin{bmatrix}1&-2&0&3\\2&2&4&7\\1&1&1&1\end{bmatrix},\quad\widetilde y=\begin{bmatrix}5\\3\\1\end{bmatrix}.\]
The last row of $H$ enforces $c_1+\cdots+c_k=1$. Row reduction of the augmented system gives
\[\operatorname{rref}[H\mid\widetilde y]=\begin{bmatrix}1&0&0&0&2\\0&1&0&- \frac{3}{2}&- \frac{3}{2}\\0&0&1&\frac{5}{2}&\frac{1}{2}\end{bmatrix}.\]
There is a free parameter; setting it to zero gives one valid choice. Other choices are also allowed, since an affine representation need not be unique for a dependent set.
One affine coefficient column is
\[c=\begin{bmatrix}2\\- \frac{3}{2}\\\frac{1}{2}\\0\end{bmatrix}.\]
Its entries sum to $1$, and substituting them in $\sum c_jv_j$ gives exactly $\begin{bmatrix}5\\3\end{bmatrix}$. Thus the requested affine combination exists.

\exercise{8.1.2}{Require the weights to sum to one}
To require both the point equation and the coefficient-sum condition, solve
\[Hc=\widetilde y,\qquad H=\begin{bmatrix}1&-1&3\\1&2&2\\1&1&1\end{bmatrix},\quad\widetilde y=\begin{bmatrix}5\\7\\1\end{bmatrix}.\]
The last row of $H$ enforces $c_1+\cdots+c_k=1$. Row reduction of the augmented system gives
\[\operatorname{rref}[H\mid\widetilde y]=\begin{bmatrix}1&0&0&-5\\0&1&0&2\\0&0&1&4\end{bmatrix}.\]
One affine coefficient column is
\[c=\begin{bmatrix}-5\\2\\4\end{bmatrix}.\]
Its entries sum to $1$, and substituting them in $\sum c_jv_j$ gives exactly $\begin{bmatrix}5\\7\end{bmatrix}$. Thus the requested affine combination exists.

\exercise{8.1.3}{Require the weights to sum to one}
To require both the point equation and the coefficient-sum condition, solve
\[Hc=\widetilde y,\qquad H=\begin{bmatrix}-3&0&4\\1&4&-2\\1&-2&6\\1&1&1\end{bmatrix},\quad\widetilde y=\begin{bmatrix}17\\1\\5\\1\end{bmatrix}.\]
The last row of $H$ enforces $c_1+\cdots+c_k=1$. Row reduction of the augmented system gives
\[\operatorname{rref}[H\mid\widetilde y]=\begin{bmatrix}1&0&0&-3\\0&1&0&2\\0&0&1&2\\0&0&0&0\end{bmatrix}.\]
One affine coefficient column is
\[c=\begin{bmatrix}-3\\2\\2\end{bmatrix}.\]
Its entries sum to $1$, and substituting them in $\sum c_jv_j$ gives exactly $\begin{bmatrix}17\\1\\5\end{bmatrix}$. Thus the requested affine combination exists.

\exercise{8.1.4}{Require the weights to sum to one}
To require both the point equation and the coefficient-sum condition, solve
\[Hc=\widetilde y,\qquad H=\begin{bmatrix}1&2&4\\2&-6&3\\0&7&1\\1&1&1\end{bmatrix},\quad\widetilde y=\begin{bmatrix}-3\\4\\-4\\1\end{bmatrix}.\]
The last row of $H$ enforces $c_1+\cdots+c_k=1$. Row reduction of the augmented system gives
\[\operatorname{rref}[H\mid\widetilde y]=\begin{bmatrix}1&0&0&\frac{13}{5}\\0&1&0&- \frac{2}{5}\\0&0&1&- \frac{6}{5}\\0&0&0&0\end{bmatrix}.\]
One affine coefficient column is
\[c=\begin{bmatrix}\frac{13}{5}\\- \frac{2}{5}\\- \frac{6}{5}\end{bmatrix}.\]
Its entries sum to $1$, and substituting them in $\sum c_jv_j$ gives exactly $\begin{bmatrix}-3\\4\\-4\end{bmatrix}$. Thus the requested affine combination exists.

\exercise{8.1.5}{Use orthogonality before checking the affine condition}
The given orthogonal basis has squared lengths $6,5,30$. Orthogonal expansion therefore gives the unique linear coefficients
\[c_1=\frac{p\cdot b_1}{6},\qquad c_2=\frac{p\cdot b_2}{5},\qquad c_3=\frac{p\cdot b_3}{30}.\]
Because the basis is linearly independent, these are the only possible coefficients. The representation is affine exactly when their sum is $1$.
\step{(a) The point $p=\begin{bmatrix}3\\8\\4\end{bmatrix}$.}
The dot products with $b_1,b_2,b_3$ are $18,\ -5,\ -30$, respectively. Dividing by $6,5,30$ gives
\[c=\begin{bmatrix}3\\-1\\-1\end{bmatrix},\qquad c_1+c_2+c_3=1.\]
Thus $p=(3)b_1+(-1)b_2+(-1)b_3$ is an affine combination.
\step{(b) The point $p=\begin{bmatrix}6\\-3\\3\end{bmatrix}$.}
The dot products with $b_1,b_2,b_3$ are $12,\ 0,\ 30$, respectively. Dividing by $6,5,30$ gives
\[c=\begin{bmatrix}2\\0\\1\end{bmatrix},\qquad c_1+c_2+c_3=3.\]
The unique linear coefficients do not sum to $1$, so no affine combination of these basis points equals $p$.
\step{(c) The point $p=\begin{bmatrix}0\\-1\\-5\end{bmatrix}$.}
The dot products with $b_1,b_2,b_3$ are $-6,\ 10,\ 0$, respectively. Dividing by $6,5,30$ gives
\[c=\begin{bmatrix}-1\\2\\0\end{bmatrix},\qquad c_1+c_2+c_3=1.\]
Thus $p=(-1)b_1+(2)b_2+(0)b_3$ is an affine combination.


\exercise{8.1.6}{Use orthogonality before checking the affine condition}
The given orthogonal basis has squared lengths $6,5,30$. Orthogonal expansion therefore gives the unique linear coefficients
\[c_1=\frac{p\cdot b_1}{6},\qquad c_2=\frac{p\cdot b_2}{5},\qquad c_3=\frac{p\cdot b_3}{30}.\]
Because the basis is linearly independent, these are the only possible coefficients. The representation is affine exactly when their sum is $1$.
\step{(a) The point $p=\begin{bmatrix}0\\-19\\-5\end{bmatrix}$.}
The dot products with $b_1,b_2,b_3$ are $-24,\ 10,\ 90$, respectively. Dividing by $6,5,30$ gives
\[c=\begin{bmatrix}-4\\2\\3\end{bmatrix},\qquad c_1+c_2+c_3=1.\]
Thus $p=(-4)b_1+(2)b_2+(3)b_3$ is an affine combination.
\step{(b) The point $p=\begin{bmatrix}\frac{3}{2}\\- \frac{13}{10}\\- \frac{1}{2}\end{bmatrix}$.}
The dot products with $b_1,b_2,b_3$ are $\frac{6}{5},\ \frac{5}{2},\ 9$, respectively. Dividing by $6,5,30$ gives
\[c=\begin{bmatrix}\frac{1}{5}\\\frac{1}{2}\\\frac{3}{10}\end{bmatrix},\qquad c_1+c_2+c_3=1.\]
Thus $p=(\frac{1}{5})b_1+(\frac{1}{2})b_2+(\frac{3}{10})b_3$ is an affine combination.
\step{(c) The point $p=\begin{bmatrix}5\\-4\\0\end{bmatrix}$.}
The dot products with $b_1,b_2,b_3$ are $6,\ 5,\ 30$, respectively. Dividing by $6,5,30$ gives
\[c=\begin{bmatrix}1\\1\\1\end{bmatrix},\qquad c_1+c_2+c_3=3.\]
The unique linear coefficients do not sum to $1$, so no affine combination of these basis points equals $p$.


\exercise{8.1.7}{Distinguish a linear span from an affine hull}
Place the three given independent vectors in the columns of
\[V=\begin{bmatrix}1&2&-1\\0&-1&2\\3&0&1\\0&4&1\end{bmatrix}.\]
A target is in $\Span S$ if $Vc=p$ is consistent; it is in $\operatorname{aff}S$ only if its coefficients also sum to $1$. Independence makes a consistent coefficient column unique.
\step{(a) Test $p_{1}$.}
\[\operatorname{rref}[V\mid p_{1}]=\begin{bmatrix}1&0&0&2\\0&1&0&1\\0&0&1&-1\\0&0&0&0\end{bmatrix}.\]
The system is consistent with $c=\begin{bmatrix}2\\1\\-1\end{bmatrix}$, so the point is in $\Span S$. The sum is $2$, so the point is not in $\operatorname{aff}S$.
\step{(b) Test $p_{2}$.}
\[\operatorname{rref}[V\mid p_{2}]=\begin{bmatrix}1&0&0&2\\0&1&0&-4\\0&0&1&3\\0&0&0&0\end{bmatrix}.\]
The system is consistent with $c=\begin{bmatrix}2\\-4\\3\end{bmatrix}$, so the point is in $\Span S$. The sum is $1$, so the point is also in $\operatorname{aff}S$.
\step{(c) Test $p_{3}$.}
\[\operatorname{rref}[V\mid p_{3}]=\begin{bmatrix}1&0&0&0\\0&1&0&0\\0&0&1&0\\0&0&0&1\end{bmatrix}.\]
The augmented last column has a pivot, so the system is inconsistent. The point is not in $\Span S$ and therefore cannot be in the smaller set $\operatorname{aff}S$.


\exercise{8.1.8}{Distinguish a linear span from an affine hull}
Place the three given independent vectors in the columns of
\[V=\begin{bmatrix}1&2&3\\0&1&0\\3&6&12\\-2&-5&-6\end{bmatrix}.\]
A target is in $\Span S$ if $Vc=p$ is consistent; it is in $\operatorname{aff}S$ only if its coefficients also sum to $1$. Independence makes a consistent coefficient column unique.
\step{(a) Test $p_{1}$.}
\[\operatorname{rref}[V\mid p_{1}]=\begin{bmatrix}1&0&0&3\\0&1&0&-1\\0&0&1&1\\0&0&0&0\end{bmatrix}.\]
The system is consistent with $c=\begin{bmatrix}3\\-1\\1\end{bmatrix}$, so the point is in $\Span S$. The sum is $3$, so the point is not in $\operatorname{aff}S$.
\step{(b) Test $p_{2}$.}
\[\operatorname{rref}[V\mid p_{2}]=\begin{bmatrix}1&0&0&0\\0&1&0&0\\0&0&1&0\\0&0&0&1\end{bmatrix}.\]
The augmented last column has a pivot, so the system is inconsistent. The point is not in $\Span S$ and therefore cannot be in the smaller set $\operatorname{aff}S$.
\step{(c) Test $p_{3}$.}
\[\operatorname{rref}[V\mid p_{3}]=\begin{bmatrix}1&0&0&-2\\0&1&0&6\\0&0&1&-3\\0&0&0&0\end{bmatrix}.\]
The system is consistent with $c=\begin{bmatrix}-2\\6\\-3\end{bmatrix}$, so the point is in $\Span S$. The sum is $1$, so the point is also in $\operatorname{aff}S$.


\exercise{8.1.9}{Represent a translated line by two points}
Write the free scalar as $t$ to avoid confusing it with a vector coordinate. Choose
\[v_1=p=\begin{bmatrix}-3\\0\end{bmatrix},\qquad v_2=p+u=\begin{bmatrix}1\\-2\end{bmatrix}.\]
Then $p+tu=(1-t)v_1+tv_2$, whose two weights sum to $1$. Conversely any affine combination of $v_1,v_2$ has precisely this form. Thus the entire solution set is $\operatorname{aff}\{v_1,v_2\}$, as required.

\exercise{8.1.10}{Represent a translated line by two points}
Write the free scalar as $t$ to avoid confusing it with a vector coordinate. Choose
\[v_1=p=\begin{bmatrix}1\\-3\\4\end{bmatrix},\qquad v_2=p+u=\begin{bmatrix}6\\-2\\2\end{bmatrix}.\]
Then $p+tu=(1-t)v_1+tv_2$, whose two weights sum to $1$. Conversely any affine combination of $v_1,v_2$ has precisely this form. Thus the entire solution set is $\operatorname{aff}\{v_1,v_2\}$, as required.

\exercise{8.1.11}{Affine closure, flats, and hull identities}
\begin{enumerate}[label=(\alph*)]
\item \textbf{True.} This is the definition of the affine hull.
\item \textbf{False.} For $S=\{e_1,e_2\}\subset\R^2$, zero is a linear combination but not an affine combination: its unique linear coefficients are $(0,0)$, whose sum is $0$, not $1$.
\item \textbf{True.} The hull is $v_1+\Span\{v_2-v_1\}$, a one-dimensional flat when $v_2\ne v_1$.
\item \textbf{False.} A flat is a translated subspace and need not contain zero. The line $\{(x,1):x\in\R\}$ is a flat but not a subspace.
\item \textbf{True.} A plane in $\R^3$ has dimension $2=3-1$, which is exactly the hyperplane dimension.
\end{enumerate}

\exercise{8.1.12}{Affine closure, flats, and hull identities}
\begin{enumerate}[label=(\alph*)]
\item \textbf{False.} The only affine combinations of one point $x$ use total coefficient $1$, so $\operatorname{aff}\{x\}=\{x\}$.
\item \textbf{True.} If a set contains its affine hull, it contains every affine combination of its points, which is the closure condition defining an affine set. Conversely an affine set contains all those combinations. Every set is contained in its own affine hull, so containment here is actually equality.
\item \textbf{True.} A one-dimensional flat is a line.
\item \textbf{False in general.} A hyperplane in $\R^n$ has dimension $n-1$. A two-dimensional flat is a hyperplane in $\R^3$, but not in $\R^4$.
\item \textbf{True.} Write the flat as $p+W$ with $W$ a subspace. If $0=p+w$ for some $w\in W$, then $p=-w\in W$, so $p+W=W$ is a subspace.
\end{enumerate}

\exercise{8.1.13}{Affine closure, flats, and hull identities}
Suppose $c_2(v_2-v_1)+c_3(v_3-v_1)=0$. Regrouping gives $(-c_2-c_3)v_1+c_2v_2+c_3v_3=0$. Since $v_1,v_2,v_3$ form a basis, all three coefficients are zero, in particular $c_2=c_3=0$. Thus the two differences are independent. Their span therefore has dimension $2$ and is a plane through zero in $\R^3$.

\exercise{8.1.14}{Affine closure, flats, and hull identities}
The affine-hull identity gives
\[
\operatorname{aff}\{v_1,v_2,v_3\}
=v_1+\Span\{v_2-v_1,v_3-v_1\}.
\]
By Exercise 13 the difference span is two-dimensional, so this is a plane. It contains $v_1$, $v_2=v_1+(v_2-v_1)$, and $v_3=v_1+(v_3-v_1)$. Conversely any plane containing those three points must contain every affine combination of them, so this is precisely the plane through them. The independence of the differences rules out a line instead.

\exercise{8.1.15}{Affine closure, flats, and hull identities}
Take solutions $x_1,\ldots,x_k$ and coefficients $c_j$ with $\sum_jc_j=1$. Matrix linearity gives
\[
A\left(\sum_jc_jx_j\right)=\sum_jc_jAx_j
=\sum_jc_jb=\left(\sum_jc_j\right)b=b.
\]
Thus every affine combination is another solution. This proves the solution set is affine. If the system is inconsistent, its solution set is empty and the closure assertion holds vacuously; if it has a solution $p$, the more concrete form is $p+\Nul A$.

\exercise{8.1.16}{Affine closure, flats, and hull identities}
For points $x_j$ satisfying $x_j\cdot v=k$ and coefficients summing to $1$,
\[
\left(\sum_jc_jx_j\right)\cdot v
=\sum_jc_j(x_j\cdot v)=\sum_jc_jk=k.
\]
Hence the set is closed under affine combinations. If $v=0$, the set is all of $\R^n$ when $k=0$, and empty otherwise; both satisfy that closure convention.

\exercise{8.1.17}{Affine closure, flats, and hull identities}
Take $S=\{(0,0,5),(1,0,5),(0,1,5)\}$. Their differences from the first point are $e_1,e_2$, which span all horizontal directions. Thus
\[
\operatorname{aff}S=(0,0,5)+\Span\{e_1,e_2\}
=\{(x_1,x_2,5):x_1,x_2\in\R\}.
\]
This is exactly the required plane, not merely a subset of it.

\exercise{8.1.18}{Affine closure, flats, and hull identities}
Choose $p=(6,0,0)$, $u=(-1,2,0)$, $v=(3,0,2)$, and the four distinct points
\[
S=\{p,p+u,p+v,p+u+v\}
=\{(6,0,0),(5,2,0),(9,0,2),(8,2,2)\}.
\]
Each satisfies $2x_1+x_2-3x_3=12$. The two directions $u,v$ are independent, so the affine hull is the plane $p+\Span\{u,v\}$. Conversely, a point on the required plane has the expression
\[
p+\frac{x_2}{2}u+\frac{x_3}{2}v
=\left(6-\frac{x_2}{2}+\frac{3x_3}{2},x_2,x_3\right),
\]
and the first coordinate equals $x_1$ by its plane equation. Thus no points of the desired plane are missing from the hull.

\exercise{8.1.19}{Affine closure, flats, and hull identities}
Take finitely many points $y_j=f(x_j)$ in $f(S)$ with $x_j\in S$, and coefficients summing to $1$. Since $S$ is affine, $x=\sum_jc_jx_j$ is in $S$. Linearity gives $\sum_jc_jy_j=f(x)$, which is in $f(S)$. Thus $f(S)$ is closed under affine combinations and is affine. If $S$ is empty, its image is empty and the same closure statement is vacuous.

\exercise{8.1.20}{Affine closure, flats, and hull identities}
For points $x_j\in S$ we know $f(x_j)\in T$. If $\sum_jc_j=1$, linearity gives $f(\sum_jc_jx_j)=\sum_jc_jf(x_j)$. Because $T$ is affine, this image belongs to $T$, hence the combined input belongs to $S$. This verifies affine closure of the preimage, including the empty-preimage case.

\exercise{8.1.21}{Affine closure, flats, and hull identities}
Every point of $\operatorname{aff}A$ is an affine combination of finitely many points of $A$. They also lie in $B$ because $A\subseteq B$. Since $B$ is affine, the combination lies in $B$. Therefore $\operatorname{aff}A\subseteq B$.

\exercise{8.1.22}{Affine closure, flats, and hull identities}
The affine set $\operatorname{aff}B$ contains $B$ and hence contains $A$. Applying Exercise 21 with that affine set gives $\operatorname{aff}A\subseteq\operatorname{aff}B$.

\exercise{8.1.23}{Affine closure, flats, and hull identities}
Both $A\subseteq A\cup B$ and $B\subseteq A\cup B$. Exercise 22 therefore gives $\operatorname{aff}A\subseteq\operatorname{aff}(A\cup B)$ and $\operatorname{aff}B\subseteq\operatorname{aff}(A\cup B)$. Taking their union proves the requested inclusion.

\exercise{8.1.24}{Affine closure, flats, and hull identities}
Take $A=\{(0,0)\}$ and $B=\{(1,0)\}$. Each singleton is its own affine hull, so their hull union is just two points. In contrast, $\operatorname{aff}(A\cup B)=\{(t,0):t\in\R\}$ is an entire line. For instance $(1/2,0)$ lies in the latter but not the former, proving strict inclusion.

\exercise{8.1.25}{Affine closure, flats, and hull identities}
Since $A\cap B\subseteq A$ and $A\cap B\subseteq B$, monotonicity of affine hull gives $\operatorname{aff}(A\cap B)\subseteq\operatorname{aff}A$ and $\operatorname{aff}(A\cap B)\subseteq\operatorname{aff}B$. Membership in both is exactly membership in their intersection, yielding the asserted inclusion.

\exercise{8.1.26}{Affine closure, flats, and hull identities}
Let $A=\{(0,0),(1,0),(0,1)\}$ and $B=\{(0,0),(2,0),(0,2)\}$. Their intersection is the singleton $\{(0,0)\}$, whose affine hull is itself. Each original set contains three noncollinear points and hence has affine hull $\R^2$. Thus $\operatorname{aff}A\cap\operatorname{aff}B=\R^2$, strictly larger than $\operatorname{aff}(A\cap B)=\{(0,0)\}$. This example does not rely on a convention for the affine hull of the empty set.
